% !TeX program = xelatex
\documentclass[12pt,a4paper]{article}
\usepackage[margin=2cm]{geometry}
\usepackage{fontspec}
\usepackage{amsmath,amssymb}
\usepackage{enumitem}
\setmainfont{Times New Roman}
\setlength{\parindent}{0pt}
\setlength{\parskip}{6pt}
\begin{document}
\begin{center}
{\Large\bfseries ĐÁP ÁN ĐỀ THI TUYỂN SINH VÀO LỚP 10 THPT}\\
{\bfseries Thành phố Hà Nội, năm học 2010--2011}\\
{\itshape Phần Đại số}
\end{center}

\textbf{Câu 1.}
\begin{enumerate}[label=\alph*)]
\item Với $x\ge0$, $x\ne9$,
\begin{align*}
P&=\frac{\sqrt{x}}{\sqrt{x}+3}
+\frac{2\sqrt{x}}{\sqrt{x}-3}
-\frac{3x+9}{(\sqrt{x}-3)(\sqrt{x}+3)}\\
&=\frac{\sqrt{x}(\sqrt{x}-3)
+2\sqrt{x}(\sqrt{x}+3)-(3x+9)}
{(\sqrt{x}-3)(\sqrt{x}+3)}\\
&=\frac{3\sqrt{x}-9}{(\sqrt{x}-3)(\sqrt{x}+3)}
=\frac{3}{\sqrt{x}+3}.
\end{align*}

\item
\[
\frac{3}{\sqrt{x}+3}=\frac13
\Longleftrightarrow \sqrt{x}+3=9
\Longleftrightarrow x=36.
\]

\item Vì $x\ge0$ nên $\sqrt{x}+3\ge3$. Do đó
\[
P=\frac3{\sqrt{x}+3}\le1.
\]
Dấu bằng xảy ra khi $x=0$. Vậy $\boxed{\max P=1}$.
\end{enumerate}

\textbf{Câu 2.}

Gọi chiều rộng của mảnh đất là $x$ mét ($0<x<13$). Chiều dài là $x+7$ mét. Theo định lý Pythagore,
\[
x^2+(x+7)^2=13^2
\Longleftrightarrow x^2+7x-60=0
\Longleftrightarrow (x-5)(x+12)=0.
\]
Đối chiếu điều kiện, $x=5$. Vậy chiều rộng là 5 m và chiều dài là 12 m.

\textbf{Câu 3.}
\begin{enumerate}[label=\alph*)]
\item Phương trình hoành độ giao điểm là
\[
-x^2=mx-1
\Longleftrightarrow x^2+mx-1=0. \tag{1}
\]
Ta có $\Delta=m^2+4>0$ với mọi $m$, nên (1) luôn có hai nghiệm phân biệt. Vì vậy $(d)$ luôn cắt $(P)$ tại hai điểm phân biệt.

\item Theo định lý Viète,
\[
x_1+x_2=-m,\qquad x_1x_2=-1.
\]
Do đó
\begin{align*}
x_1^2x_2+x_2^2x_1-x_1x_2=3
&\Longleftrightarrow x_1x_2(x_1+x_2)-x_1x_2=3\\
&\Longleftrightarrow (-1)(-m)-(-1)=3\\
&\Longleftrightarrow m=2.
\end{align*}
Vậy $\boxed{m=2}$.
\end{enumerate}

\textbf{Câu 5.}

Đặt $t=\sqrt{x^2+7}$, $t>0$. Vì $t^2=x^2+7$, phương trình trở thành
\[
t^2+4x=(x+4)t
\Longleftrightarrow t^2-(x+4)t+4x=0
\Longleftrightarrow (t-x)(t-4)=0.
\]
Nếu $t=x$ thì $\sqrt{x^2+7}=x$, điều này không thể xảy ra.

Nếu $t=4$ thì
\[
\sqrt{x^2+7}=4
\Longleftrightarrow x^2=9
\Longleftrightarrow x=\pm3.
\]
Cả hai giá trị đều thỏa mãn. Vậy
\[
\boxed{S=\{-3,3\}}.
\]
\end{document}
